\section{LoRa Adapter Rank}\label{sec:appendix-g}

<An aside comparing performance when LoRa was trained with low rank vs trained with high rank and then rank reduced. This should justify our choice of train rank and choice of rank to reduce to (they might be the same depending on results). Discuss here the benefits of low rank for training/inference but also for interp etc.>

Inspired by the fact that assistant axis found that personas can be largely described as single directions in activation space, we ran experiments investigating the required rank of the trait modifier LoRa adapters. We take our adapter LoRas and use SVD to reduce the rank from X down to 1\. See Figure X showing the trait strength in the output for each of the given ranks. Here we find that most/some of the traits can be fully modulated by LoRas as low as rank 1\. Some more complex traits require higher ranks for full behaviour, but most still remain low rank. <STILL TO DO OR PUT IN FUTURE WORK how maybe rank-1 and a small scale up is good enough to recover the full strength?> We extend this by repeating the experiment, but instead of taking the higher rank adapters and reducing the rank, we directly train the LoRa adapters at the given ranks.

\textit{Figure X showing the trait strength in the output for each of the given ranks. (downranked AND trained directly at the given rank).}

Figure X shows that you can directly train at rank 1(or 4) <or if the results say the opposite no, you need to train higher rank and then reduce>. Having lower rank is nicer for us because it makes interpretation easier, it means tweaking traits modifies less weights <or modifies them less? I think more accurately that’s not true, but modifies fewer “directions”>, meaning that it keeps the modifier adapters more independent than they otherwise might be.

